\documentclass[10pt,a4paper]{amsart}
\usepackage{amsmath,amssymb,amsthm,graphicx}
% For commutative diagrams we use Xy-pic
\usepackage[all]{xy}

% We use 2cell for 2-commutative diagrams.
\xyoption{2cell}
% Theorem environments.
%
\theoremstyle{plain}
\newtheorem{theorem}[subsection]{Theorem}
\newtheorem{proposition}[subsection]{Proposition}
\newtheorem{lemma}[subsection]{Lemma}

\theoremstyle{definition}
\newtheorem{definition}[subsection]{Definition}
\newtheorem{example}[subsection]{Example}
\newtheorem{exercise}[subsection]{Exercise}
\newtheorem{situation}[subsection]{Situation}

\theoremstyle{remark}
\newtheorem{remark}[subsection]{Remark}
\newtheorem{remarks}[subsection]{Remarks}

\numberwithin{equation}{subsection}

% Macros
%
\def\lim{\mathop{\mathrm{lim}}\nolimits}
\def\colim{\mathop{\mathrm{colim}}\nolimits}
\def\Spec{\mathop{\mathrm{Spec}}}
\def\Hom{\mathop{\mathrm{Hom}}\nolimits}
\def\Ext{\mathop{\mathrm{Ext}}\nolimits}
\def\SheafHom{\mathop{\mathcal{H}\!\mathit{om}}\nolimits}
\def\SheafExt{\mathop{\mathcal{E}\!\mathit{xt}}\nolimits}
\def\Sch{\mathit{Sch}}
\def\Mor{\mathop{\mathrm{Mor}}\nolimits}
\def\Ob{\mathop{\mathrm{Ob}}\nolimits}
\def\Sh{\mathop{\mathit{Sh}}\nolimits}
\def\NL{\mathop{N\!L}\nolimits}
\def\CH{\mathop{\mathrm{CH}}\nolimits}
\def\proetale{{pro\text{-}\acute{e}tale}}
\def\etale{{\acute{e}tale}}
\def\QCoh{\mathit{QCoh}}
\def\Ker{\mathop{\mathrm{Ker}}}
\def\Im{\mathop{\mathrm{Im}}}
\def\Coker{\mathop{\mathrm{Coker}}}
\def\Coim{\mathop{\mathrm{Coim}}}

% Boxtimes
%
\DeclareMathSymbol{\boxtimes}{\mathbin}{AMSa}{"02}

%
% Macros for moduli stacks/spaces
%
\def\QCohstack{\mathcal{QC}\!\mathit{oh}}
\def\Cohstack{\mathcal{C}\!\mathit{oh}}
\def\Spacesstack{\mathcal{S}\!\mathit{paces}}
\def\Quotfunctor{\mathrm{Quot}}
\def\Hilbfunctor{\mathrm{Hilb}}
\def\Curvesstack{\mathcal{C}\!\mathit{urves}}
\def\Polarizedstack{\mathcal{P}\!\mathit{olarized}}
\def\Complexesstack{\mathcal{C}\!\mathit{omplexes}}
% \Pic is the operator that assigns to X its picard group, usage \Pic(X)
% \Picardstack_{X/B} denotes the Picard stack of X over B
% \Picardfunctor_{X/B} denotes the Picard functor of X over B
\def\Pic{\mathop{\mathrm{Pic}}\nolimits}
\def\Picardstack{\mathcal{P}\!\mathit{ic}}
\def\Picardfunctor{\mathrm{Pic}}
\def\Deformationcategory{\mathcal{D}\!\mathit{ef}}
\usepackage{cleveref}
\crefname{originalreference}{original reference}{original references}
\Crefname{originalreference}{Original reference}{Original references}
\makeatletter
\newlabel{lemma-cd-curve-over-field}{{0F0S}{source}{}{sourcepointer-1}{}}
\newlabel{lemma-cd-curve-over-field@cref}{{[originalreference][0][]0F0S}{[1][0][]source}{}{}{}}
\makeatother
\crefname{originalreference}{original reference}{original references}
\Crefname{originalreference}{Original reference}{Original references}
\makeatletter
\newlabel{lemma-cd-limit}{{0F0R}{source}{}{sourcepointer-2}{}}
\newlabel{lemma-cd-limit@cref}{{[originalreference][0][]0F0R}{[1][0][]source}{}{}{}}
\makeatother
\title{The cohomological dimension of a field extension is bounded by the cohomological dimension of the base field plus the transcendence degree}
\author{The Stacks Project Authors}
\date{stacks-project:etale-cohomology.tex:local-snapshot-20261011; GNU FDL 1.2 or later}
\begin{document}
\maketitle
\noindent\textit{Editorial note.} All delivered mathematical text is the author's own, copied from the pinned Stacks Project chapter file etale-cohomology.tex; no statement or proof step was rewritten, reordered or removed.\par\medskip
\noindent\textit{Editorial note.} The counted proof is the author's complete argument: the tower reduction to transcendence degree one, the presentation of the field as a filtered colimit of finite type subalgebras with the resulting reduction, and the final curve bound.\par\medskip
\noindent\textit{Editorial note.} The label lemma-cd-field-extension names the counted statement, and the proof environment immediately following it in the source is delivered verbatim as its proof.\par\medskip
\noindent\textit{Editorial note.} The counted statement is carried by the author's own lemma environment in the chapter source.\par\medskip
\noindent\textit{Editorial note.} Pre-open gate: python operations/agent-authoring/tools/precard\_gate.py --source etale-cohomology.tex --label lemma-cd-field-extension --start 19672 --end 19689 --local-id 3561 returned RESULT=OK (L3561 free, occupancy FREE, public interval clear). Candidate taken from lane-s4/halfB-priority.json.\par\medskip
\noindent\textit{Editorial note.} The statement and proof are copied from the pinned chapter file; the cited inputs are kept as external pointers with their canonical Stacks tags.\par\medskip
\noindent\textit{Allowed standard prerequisites (editorial).}\begin{itemize}
\item Cohomological dimension of a filtered colimit of rings. Same chapter, the lemma bounding the cohomological dimension of a filtered colimit of algebras by the dimensions of the constituents. Supplies the reduction from L to its finite type K-subalgebras.
\item Cohomological dimension of a curve over a field. Same chapter, the lemma bounding the cohomological dimension of the function field of a curve over a field of given cohomological dimension. Supplies the bound cd(A\_i) <= 1 + cd(K) for a finite type K-subalgebra of transcendence degree one, which finishes the proof.
\end{itemize}\medskip
\ifdefined\hypertarget\hypertarget{sourcepointer-1}{}\fi
\noindent\textit{Source reference.} Original reference 0F0S: sources/stacks-etale-cohomology/pinned-original-source.tex, source lines 19651--19651. Prerequisite pointer invoked at proof line(s) 19688; it names a standard fact of the same chapter that the counted argument uses, and every step of the counted proof built on it is retained verbatim in the delivered body.\par\medskip
\ifdefined\hypertarget\hypertarget{sourcepointer-2}{}\fi
\noindent\textit{Source reference.} Original reference 0F0R: sources/stacks-etale-cohomology/pinned-original-source.tex, source lines 19633--19633. Prerequisite pointer invoked at proof line(s) 19686; it names a standard fact of the same chapter that the counted argument uses, and every step of the counted proof built on it is retained verbatim in the delivered body.\par\medskip
\begin{lemma}
\label{lemma-cd-field-extension}
Let $L/K$ be a field extension. Then we have
$\text{cd}(L) \leq \text{cd}(K) + \text{trdeg}_K(L)$.
\end{lemma}

\begin{proof}
If $\text{trdeg}_K(L) = \infty$, then this is clear.
If not then we can find a sequence of extensions
$L= L_r/L_{r - 1}/ \ldots /L_1/L_0 = K$ such that
$\text{trdeg}_{L_i}(L_{i + 1}) = 1$ and $r = \text{trdeg}_K(L)$.
Hence it suffices to prove the lemma in the case that $r = 1$.
In this case we can write $L = \colim A_i$
as a filtered colimit of its finite type $K$-subalgebras.
By Lemma \ref{lemma-cd-limit} it suffices to prove that
$\text{cd}(A_i) \leq 1 + \text{cd}(K)$. This follows
from Lemma \ref{lemma-cd-curve-over-field}.
\end{proof}

\end{document}
